\documentclass[preprint,12pt]{elsarticle}
\pdfminorversion=7

\usepackage[T1]{fontenc}
\usepackage{lmodern}
\usepackage{booktabs}
\usepackage{tabularx}
\usepackage{array}
\usepackage{graphicx}
\usepackage{amsmath}
\usepackage[hidelinks]{hyperref}
\usepackage{xcolor}
\usepackage{url}
\setlength{\emergencystretch}{2em}

\journal{Expert Systems with Applications}

\newcommand{\sourcefigure}[3]{%
\begin{figure}[t]
\centering
\includegraphics[width=0.94\linewidth]{#2}
\caption{#3}
\label{fig:#1}
\end{figure}
}

\begin{document}

\begin{frontmatter}

\title{Protocol-Aware Expert-System Validation for COVID-19 Respiratory-Audio Screening Under Temporal Drift, External Dataset Shift, and Metadata Confounding}

\author[inst1]{Author names withheld for review}
\address[inst1]{Institutional affiliation withheld for review}

\begin{abstract}
Respiratory-audio screening has been proposed as a low-cost decision-support route for COVID-19, yet high internal accuracy can arise from participant leakage, collection-period effects, dataset identity, symptom correlations, and protocol metadata rather than portable disease acoustics. This paper presents a protocol-aware expert-system validation framework for COVID-19 respiratory-audio models. The framework couples a broad acoustic model bank with participant-disjoint testing, time-aware validation, strict chronological testing, cough-only external transfer, and metadata-only confounding controls, so each prediction result is interpreted with its validation tier and claim boundary. Coswara cough, breath, and speech recordings support source-domain multimodal and temporal analyses. COUGHVID is used only as an independent cough-only transfer endpoint and is not treated as evidence for multimodal fusion. In Coswara, the strongest ComParE2016+IS10 fusion result achieved AUROC 0.897 and AUPRC 0.863. Performance decreased under a time-stratified participant split (AUROC 0.849) and under strict early-to-late temporal validation (AUROC 0.698). In model-matched cough-only comparisons, Coswara internal AUROC was 0.849-0.868, whereas COUGHVID transfer AUROC was 0.523-0.543. Metadata-only controls reached AUROC 0.964, and symptoms-only metadata reached AUROC 0.932. Inverse-propensity sensitivity reduced quality-weighted audio AUROC from 0.879 to 0.780-0.807 after measured adjustment. These results show that an internally strong respiratory-audio expert system can still be fragile under temporal shift, external dataset shift, and measured confounding.
\end{abstract}

\begin{keyword}
COVID-19 audio \sep expert systems \sep clinical decision support \sep multimodal fusion \sep temporal validation \sep cross-dataset validation \sep confounding
\end{keyword}

\end{frontmatter}

\section{Introduction}

Respiratory-audio screening systems are attractive because cough, breath, and speech can be collected with commodity devices and analyzed by automated models. In an expert-system setting, however, the relevant question is not whether a model can score highly on a convenient internal split. The relevant question is whether the decision rule remains credible when the data collection period changes, when the validation data come from a different respiratory-audio corpus, and when non-audio information such as symptoms, demographics, recording protocol, and geography can predict the label.

COVID-19 audio datasets make this distinction important. Coswara provides cough, breath, speech, and metadata for remote SARS-CoV-2 screening \citep{bhattacharya2023coswara}. COUGHVID provides a large crowdsourced cough corpus for cough-analysis research \citep{orlandic2021coughvid}. COVID-19 Sounds and related longitudinal work show that respiratory audio can be collected at scale and over time \citep{xia2021covid,dang2022longitudinal,han2022sounds}. Yet observational collection also creates a setting in which disease status can be entangled with symptoms, recording date, collection campaign, country, device, participant behavior, and inclusion criteria. A strong internal classifier may therefore be a strong detector of the study protocol rather than a robust screening aid. This concern is consistent with recent evidence that audio classifiers may lose much of their apparent advantage over symptom checkers under realistic confounding control \citep{coppock2024audio}.

This manuscript is written for Expert Systems with Applications as a decision-support reliability study. The expert-system object is the validation framework: an auditable pipeline that links feature extraction, model selection, fusion, thresholding, temporal testing, external transfer, and metadata controls. The framework is designed to prevent a common failure mode in applied machine learning, where an internally optimized model is presented as a deployable system before its protocol dependence has been tested. The output of the system is therefore not just a probability score; it is a probability score accompanied by the validation protocol under which it was produced and the claim level that the protocol can support.

The study addresses four research questions:
\begin{enumerate}
\item How strong is source-domain respiratory-audio performance when cough, breath, and speech branches are evaluated under participant-disjoint Coswara validation?
\item How much does performance change when the validation protocol adds collection-time stratification and strict early-to-late chronological testing?
\item Does a Coswara-trained cough model transfer to COUGHVID, and what does that result say about external dataset shift rather than multimodal fusion?
\item How much of the apparent prediction problem can be solved by measured symptoms, demographic variables, and recording-protocol metadata alone?
\end{enumerate}

The contributions are as follows:
\begin{enumerate}
\item A protocol-aware expert-system evaluation framework for COVID-19 respiratory audio, in which model development is explicitly separated from temporal, external, and metadata-confounding audits.
\item A reproducible model bank spanning MFCC-style baselines, spectrogram CNNs, WavLM self-supervised representations, particle-swarm feature search, OpenSMILE ComParE2016+IS10 descriptors, support vector machines, gradient-boosted trees, and late fusion.
\item A modality-scoped validation design: Coswara supports multimodal source-domain and temporal conclusions, while COUGHVID supports only cough-only external transfer.
\item A protocol-aware comparison with prior COVID-audio and ESWA DNDT/DNDF evidence, distinguishing close internal comparisons from stricter temporal and external stress tests.
\item A metadata audit showing that symptoms and demographic/protocol variables can predict labels strongly enough to constrain any causal interpretation of audio performance.
\end{enumerate}

The remainder of the paper reviews related work, describes the datasets and endpoint scope, presents the expert-system workflow and validation ladder, reports the main results, compares them to prior work by protocol, and discusses the reliability implications.

\sourcefigure{pipeline-diagram}{../common_figures/fig_pipeline.pdf}{Expert-system workflow for COVID-19 respiratory-audio reliability evaluation. Coswara cough, breath, speech, and metadata enter the source-domain branches; feature extraction and model selection feed late fusion; the final evidence is organized by participant-disjoint validation, time-stratified validation, strict early-to-late temporal testing, cough-only COUGHVID transfer, and metadata-only confounding controls.}

\section{Related Work and Research Gap}

\subsection{COVID-19 respiratory-audio datasets}

Coswara is a multimodal respiratory-audio and symptom dataset with cough, breathing, vowel, counting, and metadata fields \citep{bhattacharya2023coswara}. It is therefore suitable for asking whether multiple respiratory modalities can improve source-domain discrimination under participant-disjoint validation. COUGHVID is a separate crowdsourced cough corpus with different collection and labeling characteristics \citep{orlandic2021coughvid}. Its role in this paper is external transfer, not multimodal validation. COVID-19 Sounds and later analyses broaden the field by demonstrating large-scale collection and longitudinal respiratory-audio modeling \citep{xia2021covid,dang2022longitudinal,han2022sounds}. These datasets establish the application opportunity, but they also highlight why protocol-aware evaluation is necessary: each corpus reflects a different recruitment process, recording interface, label mechanism, and sampling period.

\subsection{Feature, self-supervised, and transformer methods}

COVID-audio modeling has used both handcrafted and learned representations. Standard acoustic baselines include MFCCs and paralinguistic descriptors. OpenSMILE provides reproducible acoustic feature extraction \citep{eyben2010opensmile}, while the IS10 and ComParE feature sets are established high-dimensional paralinguistic representations \citep{schuller2010is10,schuller2016compare}. Deep and transfer-learning approaches have also reported strong internal results for cough, breath, and speech \citep{pahar2021global,pahar2022transfer}. Spectrogram transformer architectures such as HST emphasize representation learning and attention over respiratory sounds \citep{aytekin2024hst}. Self-supervised speech models such as WavLM provide another route to learned audio embeddings \citep{chen2022wavlm}. These methods motivate a broad model bank, but representation strength alone does not answer whether a model survives participant, temporal, or dataset shift.

\subsection{Decision trees, forests, and expert systems}

Tree and forest methods remain central to many applied expert systems because they provide strong tabular performance, support feature-ranking workflows, and integrate naturally with engineered descriptors. Random forests \citep{breiman2001random}, support vector machines \citep{cortes1995svm}, LightGBM \citep{ke2017lightgbm}, XGBoost \citep{chen2016xgboost}, and CatBoost \citep{prokhorenkova2018catboost} are therefore included as practical decision modules. SMOTE \citep{chawla2002smote} and particle swarm optimization \citep{kennedy1995pso} are evaluated as class-imbalance and feature-search tools. The most relevant ESWA benchmark is the DNDT/DNDF study, which reports stronger internal AUCs for Coswara and COUGHVID cough classification, while also reporting limited bidirectional cross-dataset transfer between Coswara and COUGHVID \citep{islam2026dndf}. The present study does not dispute the internal strength of that architecture; it asks a different expert-system question by placing internal, temporal, external, and metadata-control evidence in a single validation ladder.

\subsection{Temporal shift, external validation, and confounding}

Temporal drift and shortcut learning are central risks for respiratory-audio decision support. Drift-adaptive COVID cough work reports performance degradation between development and post-development periods and motivates explicit monitoring of temporal change \citep{ganitidis2025drift}. Confounding critiques show that symptom and participant factors can explain much of the label structure in vocal-audio COVID screening \citep{coppock2024audio}. Symptom-assisted models can improve prediction by deliberately adding clinical symptom information \citep{lin2025sympcoughnet}; in this study, symptom metadata play a different role, serving as a confounding audit rather than as a desired diagnostic input. The research gap is therefore not the absence of another classifier. It is the need for an ESWA-style expert-system evaluation design that makes protocol dependence visible before a respiratory-audio model is treated as reliable decision support.

\begin{table}[t]
\centering
\caption{Thematic comparison to prior COVID respiratory-audio studies. The comparison is organized by evidence role because raw metrics are not interchangeable across datasets, modalities, split units, and external-validation assumptions.}
\label{tab:prior-work}
\scriptsize
\setlength{\tabcolsep}{2.5pt}
\begin{tabularx}{\linewidth}{p{0.16\linewidth} p{0.18\linewidth} p{0.21\linewidth} X}
\toprule
Study & Dataset/modality & Main protocol emphasis & Relation to this manuscript \\
\midrule
Coswara \citep{bhattacharya2023coswara} & Cough, breath, speech, symptoms, metadata & Dataset construction for remote screening & Source corpus for multimodal participant-disjoint, temporal, and metadata-confounding analyses. \\
COUGHVID \citep{orlandic2021coughvid} & Crowdsourced cough & Large cough corpus and cough-analysis tools & Independent cough-only transfer endpoint; not evidence for multimodal fusion. \\
COVID-19 Sounds \citep{xia2021covid,han2022sounds} & Cough, breath, voice & Large-scale digital respiratory screening and realistic audio testing & Establishes field context for scalable audio collection and cautious performance interpretation. \\
Drift work \citep{dang2022longitudinal,ganitidis2025drift} & Respiratory audio over time & Sequential modeling and drift-adaptive evaluation & Motivates collection-time stratification and strict chronological testing. \\
Pahar et al. \citep{pahar2021global,pahar2022transfer} & Cough, breath, speech & Smartphone cough and transfer-learning respiratory-audio models & Motivates multimodal modeling; our emphasis is reliability under stricter validation. \\
HST \citep{aytekin2024hst} & Respiratory sounds & Spectrogram transformer architecture & Architecture-forward reference; our paper is validation-forward. \\
SympCoughNet \citep{lin2025sympcoughnet} & Cough plus symptoms & Symptom-assisted prediction & Uses symptoms as model input; this paper uses symptoms as a confounding control. \\
DNDT/DNDF ESWA \citep{islam2026dndf} & Multiple cough datasets & Deep decision tree/forest, internal and cross-dataset cough evaluation & Reports stronger internal metrics and limited Coswara--COUGHVID transfer; we compare by protocol rather than raw leaderboard rank. \\
Confounding critique \citep{coppock2024audio} & Vocal audio and symptoms & Matched and unmatched confounding analysis & Motivates symptoms-only and demographic/protocol metadata baselines. \\
This study & Coswara and COUGHVID & Participant, time, temporal, external, and metadata-control validation ladder & Shows internal fusion AUROC 0.897, temporal AUROC 0.698, matched cough external AUROC 0.523-0.543, and metadata-only AUROC 0.964. \\
\bottomrule
\end{tabularx}
\end{table}

\section{Materials and Endpoint Scope}

\subsection{Coswara source-domain data}

Coswara is the source dataset for model development and internal validation. It contains respiratory and speech recordings together with symptom, demographic, and protocol metadata \citep{bhattacharya2023coswara}. Participant identity is treated as the leakage-control unit whenever participant identifiers are available. The source-domain audio analyses use cough, breath, and speech modalities according to endpoint availability. The final existing participant-split result uses cough+speech stacked logistic fusion; the time-stratified participant split uses cough+breath+speech uniform fusion; the strict early-to-late temporal audio endpoint is selected as a breath top-4 validation ensemble.

\subsection{COUGHVID external-transfer data}

COUGHVID is used only as an external cough-transfer test \citep{orlandic2021coughvid}. No COUGHVID recordings are pooled into Coswara training or validation. The final external-transfer endpoint contains 8331 COUGHVID cough samples. Because COUGHVID differs from Coswara in collection mechanism, label construction, and class prevalence, this endpoint is interpreted as a dataset-shift stress test. It does not validate Coswara multimodal fusion, because breath and speech modalities are not present in the external endpoint.

\begin{table}[t]
\centering
\caption{Dataset and validation-role summary for the ESWA expert-system framework. Endpoint sample sizes are taken from the reproducible result tables; source-dataset scale is cited from the dataset papers.}
\label{tab:dataset-protocol}
\scriptsize
\setlength{\tabcolsep}{2.5pt}
\begin{tabularx}{\linewidth}{p{0.16\linewidth} p{0.19\linewidth} p{0.15\linewidth} p{0.22\linewidth} X}
\toprule
Protocol row & Modalities used in this manuscript & Endpoint $n$ & Validation role & Scope of the claim \\
\midrule
Coswara source dataset & Cough, breath, speech, metadata & Dataset scale reported by Coswara & Source-domain model development and validation & Supports multimodal and temporal analysis within the Coswara collection process. \\
Existing participant split & Cough+speech & 314 & Participant-disjoint internal test & Best selected source-domain fusion result; not external validation. \\
Time-stratified participant split & Cough, breath, speech & 431 & Participant-disjoint validation with collection-time stratification & Tests whether source-domain performance remains stable after reducing simple calendar imbalance. \\
Strict early-to-late split & Breath for selected audio endpoint; metadata for ablation & 411 audio; 3816 metadata & Chronological future-period stress test & Tests temporal reliability inside Coswara; not a separate dataset. \\
COUGHVID external transfer & Cough only & 8331 & Direct Coswara-to-COUGHVID transfer & Tests cough-only dataset shift; does not evaluate breath, speech, or multimodal fusion. \\
\bottomrule
\end{tabularx}
\end{table}

\section{Protocol-Aware Expert-System Framework}

\subsection{System design}

Figure~\ref{fig:pipeline-diagram} defines the framework. The decision-support pipeline has four functional layers. The data layer separates Coswara source-domain evidence from COUGHVID external cough evidence. The representation layer extracts multiple audio feature families so that the validation results are not tied to one handcrafted or deep representation. The decision layer trains classifier families and selects thresholds using validation data. The audit layer subjects the selected predictions to increasingly severe validation protocols and compares audio performance with metadata-only controls.

\sourcefigure{evaluation-protocol-ladder}{../common_figures/fig_evaluation_ladder.pdf}{Validation ladder used to evaluate respiratory-audio expert-system reliability. The first two tiers are participant-disjoint Coswara tests, the third tier is strict early-to-late Coswara temporal testing, and the fourth tier is cough-only COUGHVID transfer. Increasing protocol severity changes the evidence target, so metrics should not be read as a single leaderboard.}

\subsection{Feature extraction and model families}

The feature bank intentionally spans simple, engineered, and learned representations. MFCC-style descriptors and spectrogram CNNs serve as conventional audio baselines. WavLM embeddings test whether self-supervised speech representations transfer useful information to respiratory-audio tasks \citep{chen2022wavlm}. Particle swarm optimization is evaluated as a wrapper-style feature-search strategy \citep{kennedy1995pso}. The final selected branch uses OpenSMILE ComParE2016+IS10 descriptors, combining standardized paralinguistic features from two established sets \citep{eyben2010opensmile,schuller2010is10,schuller2016compare}. The final top-800 feature strategy was selected from the Coswara source-domain training workflow and frozen before COUGHVID transfer; COUGHVID labels were not used for feature selection. ComParE+IS10 was used for the external-transfer endpoint because it is a standardized handcrafted representation aligned with much of the respiratory-audio baseline literature and can be applied to the target corpus without target-domain tuning. In the chronological analysis, this is a frozen source-domain representation rather than a fully prospective feature-selection protocol re-estimated only inside the early-period cohort.

The deep branches are reported as source-domain development evidence, not as external-transfer evidence. The practical WavLM and CNN-BiGRU rows are isolated single-stream checks, whereas the 0.897 AUROC endpoint is a validation-selected Coswara multimodal fusion row. This difference explains why the final fusion endpoint can exceed the single-modality deep rows without implying that WavLM or CNN-BiGRU was externally validated on COUGHVID.

The classifier bank includes logistic regression, random forests, support vector machines, LightGBM, XGBoost, CatBoost, and validation-selected top-4 ensembles. These families are deliberately practical for expert-system construction: they support tabular acoustic features, tolerate moderate sample sizes, and allow validation-driven model selection. SMOTE is used in selected boosted-tree configurations to address class imbalance \citep{chawla2002smote}.

\subsection{Fusion and thresholding}

Fusion is performed only after branch-level predictions have been generated. Depending on the validation-selected row, multimodal predictions are combined by uniform averaging, validation-weighted averaging, or stacked logistic fusion. This late-fusion design keeps modality evidence explicit: a cough+speech fusion model is a Coswara source-domain endpoint, while a COUGHVID result remains cough-only. Thresholds for sensitivity, specificity, F1, and balanced accuracy are selected on validation predictions, with balanced accuracy used as the primary threshold-selection criterion unless a specific branch defines otherwise. AUROC and AUPRC are reported separately from thresholded metrics because they evaluate ranking behavior and are not determined by a single operating point.

\subsection{Validation ladder and metrics}

The validation ladder has four tiers. Tier 1 evaluates source-domain participant-disjoint performance on Coswara. Tier 2 repeats participant-disjoint validation with collection-time stratification. Tier 3 trains and validates on earlier Coswara recordings and tests on later recordings. Tier 4 transfers a Coswara-trained cough model directly to COUGHVID. Metrics are reported as AUROC, AUPRC, balanced accuracy, F1, sensitivity, and specificity where available. AUPRC is interpreted together with prevalence and modality scope because it is sensitive to the positive-class rate.

\subsection{Expert-system claim control}

The framework assigns each result to an admissible claim level (Table~\ref{tab:claim-control}). This claim-control layer is the expert-system contribution: it prevents a source-domain result, a temporal stress-test result, and an external-transfer result from being interpreted as equivalent evidence. It also prevents symptom or protocol metadata from being silently treated as respiratory acoustic evidence.

\begin{table}[t]
\centering
\caption{Claim-control layer used by the expert-system framework. The output of a model is interpreted jointly with the validation tier and modality scope, so the strongest numerical row does not automatically become the strongest deployment claim.}
\label{tab:claim-control}
\scriptsize
\setlength{\tabcolsep}{2.5pt}
\begin{tabularx}{\linewidth}{p{0.22\linewidth} p{0.24\linewidth} p{0.18\linewidth} X}
\toprule
Evidence source & Decision-support question & Observed result & Permitted claim \\
\midrule
Coswara participant-disjoint fusion & Can source-domain multimodal branches discriminate labels without participant overlap? & AUROC 0.897 & Source-domain feasibility only. \\
Coswara time-aware and strict temporal validation & Does performance survive calendar structure and future-period shift? & AUROC 0.849 and 0.698 & Temporal robustness is incomplete and must be reported. \\
COUGHVID transfer & Does a Coswara-trained cough branch transfer to an independent cough corpus? & AUROC 0.543, AUPRC 0.040 & External cough portability is weak; multimodal external validation is not tested. \\
Metadata-only controls & Can non-audio fields predict labels without respiratory audio? & Full metadata AUROC 0.964 & Shortcut-learning and measured-confounding risk must constrain interpretation. \\
\bottomrule
\end{tabularx}
\end{table}

\subsection{Metadata-confounding audit}

The audit trains metadata-only models using full safe metadata, symptoms-only metadata, and demographic/protocol-only metadata. These controls answer whether measured non-audio variables can predict the label without respiratory audio. A strict temporal metadata ablation also compares full safe metadata with and without recording month. The support analyses identify recording year, recording month, country, age, and duration as important demographic/protocol drivers, and residual imbalance remains after inverse-probability weighting for variables such as recording year and country. These findings are used to bound interpretation, not to claim that every audio model uses the same shortcuts.

\section{Results}

\subsection{Validation severity changes the apparent reliability of the expert system}

The strongest source-domain result is the existing participant split: cough+speech stacked logistic fusion achieved AUROC 0.897, AUPRC 0.863, balanced accuracy 0.825, and F1 0.752 (Table~\ref{tab:final-protocols}). This is the best evidence that the acoustic model bank can learn source-domain Coswara signal under participant separation.

That result is not stable across stricter validation tiers. The time-stratified participant split reached AUROC 0.849 and AUPRC 0.783. Strict early-to-late temporal validation reached AUROC 0.698, despite AUPRC 0.896. The AUROC/AUPRC divergence is expected to be possible under changed class prevalence and temporal composition, and it is why the framework reports ranking, precision-recall, balanced accuracy, sensitivity, and specificity together rather than selecting one favorable metric.

\begin{table}[t]
\centering
\caption{Final selected ComParE2016+IS10 validation results across protocol tiers. The modality column defines the evidence scope: Coswara rows support source-domain and temporal analysis, while the COUGHVID row is cough-only external transfer.}
\label{tab:final-protocols}
\scriptsize
\setlength{\tabcolsep}{2.5pt}
\resizebox{\linewidth}{!}{%
\begin{tabular}{p{0.22\linewidth} p{0.17\linewidth} p{0.24\linewidth} r r r r r r r}
\toprule
Protocol & Modality & Selected model & $n$ & AUROC & AUPRC & BA & F1 & Sens. & Spec. \\
\midrule
Existing participant split & Cough+speech & Stacked logistic fusion & 314 & 0.897 & 0.863 & 0.825 & 0.752 & 0.833 & 0.816 \\
Time-stratified participant split & Cough, breath, speech & Uniform mean fusion & 431 & 0.849 & 0.783 & 0.783 & 0.705 & 0.710 & 0.857 \\
Strict early-to-late temporal split & Breath & Top-4 validation ensemble & 411 & 0.698 & 0.896 & 0.656 & 0.751 & 0.646 & 0.667 \\
COUGHVID external transfer & Cough & LightGBM & 8331 & 0.543 & 0.040 & 0.510 & 0.058 & 0.084 & 0.936 \\
\bottomrule
\end{tabular}%
}
\end{table}

\sourcefigure{temporal-degradation-plot}{../common_figures/fig_validation_degradation.pdf}{Protocol-severity summary for selected final ComParE2016+IS10 endpoints. The selected endpoints differ in modality and fusion rule, so the plot is a claim-severity summary rather than a fixed-model ablation. Table~\ref{tab:fixed-cough-ladder} reports the matched cough-only model-family ladder.}

\subsection{External transfer is a cough-only dataset-shift stress test}

Direct COUGHVID transfer is limited for all final cough-only ComParE2016+IS10 models (Table~\ref{tab:external-transfer}). The best selected external result is LightGBM with AUROC 0.543, AUPRC 0.040, balanced accuracy 0.510, sensitivity 0.084, and specificity 0.936. This endpoint should be interpreted by modality scope: COUGHVID does not supply the breath and speech branches used by Coswara fusion. It asks the stricter cough-domain question of whether a Coswara cough decision rule survives an independent cough corpus. The experiment deliberately tests a frozen ComParE+IS10 handcrafted pipeline, not every possible deep representation. It is therefore a standardized cough-transfer stress test for a common paralinguistic-feature paradigm under maximum dataset shift. Under the selected operating point, the evidence indicates weak external cough portability rather than a useful screening rule.

The selected validation ladder is useful for claim control, but it is not a fixed-model ablation because the selected modality and fusion rule change across tiers. The fair comparison is therefore cough-only internal performance against cough-only external transfer. Table~\ref{tab:cough-matched-transfer} shows that Coswara cough-only internal AUROC was 0.849-0.868 for the matched model families, while the corresponding COUGHVID AUROC was 0.523-0.543. Table~\ref{tab:fixed-cough-ladder} extends this matched comparison across the full validation ladder while holding modality fixed to cough and model family fixed. This prevents the misleading interpretation that the observed external weakness is simply caused by comparing multimodal Coswara fusion with a cough-only target; the matched cough branch itself shows poor portability.

\begin{table}[t]
\centering
\caption{COUGHVID external-transfer results for Coswara-trained cough-only ComParE2016+IS10 models. Low AUPRC and low sensitivity at the validation-selected threshold indicate that these cough models do not operate as useful external screening rules on COUGHVID.}
\label{tab:external-transfer}
\scriptsize
\setlength{\tabcolsep}{2.5pt}
\resizebox{\linewidth}{!}{%
\begin{tabular}{p{0.24\linewidth} p{0.24\linewidth} r r r r r r r}
\toprule
Model & Feature strategy & $n$ & AUROC & AUPRC & BA & F1 & Sens. & Spec. \\
\midrule
LightGBM with SMOTE & ComParE+IS10 top-800 LightGBM selection & 8331 & 0.543 & 0.040 & 0.510 & 0.058 & 0.084 & 0.936 \\
SVM RBF & ComParE+IS10 top-800 LightGBM selection & 8331 & 0.523 & 0.037 & 0.500 & 0.058 & 0.193 & 0.807 \\
CatBoost with SMOTE & ComParE+IS10 top-800 LightGBM selection & 8331 & 0.531 & 0.040 & 0.506 & 0.059 & 0.144 & 0.869 \\
XGBoost with SMOTE & ComParE+IS10 top-800 LightGBM selection & 8331 & 0.532 & 0.039 & 0.505 & 0.063 & 0.267 & 0.743 \\
\bottomrule
\end{tabular}%
}
\end{table}

\begin{table}[t]
\centering
\caption{Model-matched cough-only internal-to-external comparison. Coswara rows are participant-disjoint cough-only internal endpoints with ComParE+IS10 top-800 features. COUGHVID rows apply the corresponding Coswara-trained cough model to the independent cough-only target.}
\label{tab:cough-matched-transfer}
\small
\begin{tabular}{lcccccc}
\toprule
Model & \multicolumn{3}{c}{Coswara cough-only} & \multicolumn{3}{c}{COUGHVID cough-only} \\
\cmidrule(lr){2-4}\cmidrule(lr){5-7}
 & AUROC & AUPRC & BA & AUROC & AUPRC & BA \\
\midrule
LightGBM SMOTE & 0.853 & 0.797 & 0.782 & 0.543 & 0.040 & 0.510 \\
SVC RBF & 0.868 & 0.812 & 0.795 & 0.523 & 0.037 & 0.500 \\
CatBoost SMOTE & 0.849 & 0.788 & 0.776 & 0.531 & 0.040 & 0.506 \\
XGBoost SMOTE & 0.850 & 0.780 & 0.776 & 0.532 & 0.039 & 0.505 \\
\bottomrule
\end{tabular}
\end{table}

\begin{table}[t]
\centering
\caption{Fixed cough-only model-family ladder. Each row keeps the modality fixed to cough and follows the same model family across validation tiers. Models are retrained within the corresponding source-domain protocol; values are AUROC.}
\label{tab:fixed-cough-ladder}
\small
\resizebox{\linewidth}{!}{%
\begin{tabular}{lcccc}
\toprule
Model & Existing & Time-stratified & Strict temporal & COUGHVID \\
\midrule
LightGBM SMOTE & 0.853 & 0.830 & 0.604 & 0.543 \\
SVC RBF & 0.868 & 0.802 & 0.561 & 0.523 \\
CatBoost SMOTE & 0.849 & 0.826 & 0.614 & 0.531 \\
XGBoost SMOTE & 0.850 & 0.828 & 0.599 & 0.532 \\
\bottomrule
\end{tabular}
}
\end{table}

\subsection{Model breadth supports the audit but does not replace it}

Table~\ref{tab:model-bank} shows representative model-family rows from the reproducible outputs. Spectrogram CNN rows reached AUROC 0.730 for cough and 0.710 for breath. WavLM rows reached AUROC 0.790 for shallow breath and 0.790 for shallow cough. The ComParE2016+IS10 fusion row reached AUROC 0.897 under the existing participant split. These results justify using a broad model bank, but they also show why the manuscript should not be read as a simple model leaderboard. Rows differ in representation, modality, and validation role unless explicitly matched. In particular, the WavLM and CNN rows are source-domain single-stream checks; they are not COUGHVID transfer rows and are not direct ablations of the multimodal fusion endpoint.

\begin{table}[t]
\centering
\caption{Representative model-family rows from the reproducible analysis outputs. These rows show the breadth of the expert-system model bank; they are not direct head-to-head claims unless the protocol and endpoint match.}
\label{tab:model-bank}
\scriptsize
\setlength{\tabcolsep}{2.5pt}
\resizebox{\linewidth}{!}{%
\begin{tabular}{p{0.20\linewidth} p{0.20\linewidth} p{0.17\linewidth} r r r r r}
\toprule
Family & Representative row & Modality & $n$ & AUROC & AUPRC & BA & F1 \\
\midrule
CNN/spectrogram & CNN-BiGRU & Cough & 636 & 0.730 & 0.556 & 0.686 & 0.586 \\
CNN/spectrogram & CNN-BiGRU & Breath & 636 & 0.710 & 0.590 & 0.649 & 0.550 \\
SSL transformer & WavLM base plus & Shallow breath & 296 & 0.790 & 0.654 & 0.731 & 0.627 \\
SSL transformer & WavLM base plus & Shallow cough & 310 & 0.790 & 0.731 & 0.713 & 0.614 \\
ComParE+IS10 fusion & Cough+speech stacked logistic & Multimodal & 314 & 0.897 & 0.863 & 0.825 & 0.752 \\
\bottomrule
\end{tabular}%
}
\end{table}

\sourcefigure{multimodal-fusion-schematic}{../common_figures/fig_multimodal_fusion.pdf}{Endpoint modality scope for the fusion experiments. Cough, breath, and speech fusion is evaluated only inside Coswara; COUGHVID is a cough-only transfer endpoint and therefore cannot validate multimodal fusion.}

\subsection{Metadata-only controls expose shortcut-learning risk}

Metadata-only models are strong under the participant split (Table~\ref{tab:metadata}). Full safe metadata reached AUROC 0.964 and AUPRC 0.928. Symptoms-only metadata reached AUROC 0.932 and AUPRC 0.898. Demographic/protocol-only metadata reached AUROC 0.914. These values are higher than the final internal audio/fusion AUROC of 0.897, which means that non-audio variables can explain much of the label structure in this observational setting.

The temporal metadata ablation shows that recording month is not a simple nuisance variable. Full safe metadata reached AUROC 0.531 under strict early-to-late validation, while removing recording month increased AUROC to 0.779. This pattern indicates that calendar variables can encode collection-period structure that does not generalize chronologically. The correct conclusion is not that audio contains no signal; it is that audio-only and fusion claims must be bounded by measured confounding and temporal-protocol risk.

\subsection{Inverse-propensity sensitivity qualifies the residual audio signal}

Inverse-propensity weighting was used as a measured-confounding sensitivity analysis rather than a causal proof (Table~\ref{tab:ipw}). The unweighted quality-weighted audio fusion row reached AUROC 0.879 and AUPRC 0.832 on 318 held-out samples. With cap 2 and 95th-percentile truncation, AUROC remained 0.807 with effective sample size 238.7. With cap 10 and 99th-percentile truncation, AUROC was 0.780, AUPRC was 0.537, effective sample size was 130.4, and maximum absolute residual standardized mean difference decreased to 0.382. The drop from 0.879 to 0.780-0.807 supports a qualified interpretation: measured adjustment leaves above-chance signal, but a substantial portion of naive source-domain performance is compatible with observational selection and residual imbalance.

\begin{table}[t]
\centering
\caption{Inverse-propensity sensitivity analysis for quality-weighted audio fusion. ESS denotes effective sample size and SMD denotes standardized mean difference.}
\label{tab:ipw}
\scriptsize
\setlength{\tabcolsep}{3pt}
\resizebox{\linewidth}{!}{%
\begin{tabular}{lrrrrrr}
\toprule
Weighting setting & AUROC & AUPRC & BA & ESS & Mean abs. SMD & Max abs. SMD \\
\midrule
Unweighted & 0.879 & 0.832 & 0.804 & 318.0 & 0.224 & 1.384 \\
IPW cap 2, q=0.95 & 0.807 & 0.624 & 0.742 & 238.7 & 0.149 & 0.724 \\
IPW cap 5, q=0.95 & 0.793 & 0.573 & 0.731 & 173.2 & 0.129 & 0.523 \\
IPW cap 10, q=0.99 & 0.780 & 0.537 & 0.721 & 130.4 & 0.118 & 0.382 \\
\bottomrule
\end{tabular}%
}
\end{table}

\begin{table}[t]
\centering
\caption{Metadata-confounding evidence from participant-split and strict temporal analyses. Metadata-only performance is reported as an audit of shortcut risk, not as a proposed diagnostic input for the respiratory-audio expert system.}
\label{tab:metadata}
\scriptsize
\setlength{\tabcolsep}{2.5pt}
\resizebox{\linewidth}{!}{%
\begin{tabular}{p{0.31\linewidth} p{0.22\linewidth} r r r r r}
\toprule
Metadata configuration & Protocol & $n$ & AUROC & AUPRC & BA & F1 \\
\midrule
Full safe metadata & Existing participant split & 2862 & 0.964 & 0.928 & 0.890 & 0.849 \\
Symptoms only & Existing participant split & 2862 & 0.932 & 0.898 & 0.912 & 0.876 \\
Demographic/protocol only & Existing participant split & 2862 & 0.914 & 0.737 & 0.827 & 0.751 \\
Full safe metadata & Strict early-to-late & 3816 & 0.531 & 0.837 & 0.492 & 0.624 \\
Full safe metadata without month & Strict early-to-late & 3816 & 0.779 & 0.935 & 0.723 & 0.819 \\
\bottomrule
\end{tabular}%
}
\end{table}

\sourcefigure{confounding-evidence-diagram}{../common_figures/fig_metadata_confounding.pdf}{Confounding evidence for the respiratory-audio expert-system audit. Metadata-only baselines outperform internal audio/fusion AUROC, and the strict temporal month ablation shows that collection-period variables can change chronological generalization.}

\subsection{Protocol-aware comparison with prior results}

Table~\ref{tab:metric-comparison} separates close comparisons from apples-to-oranges comparisons. The closest paper-style internal comparison to the DNDT/DNDF ESWA result is an internal cough comparison: DNDF reports Coswara cough AUC 0.92, while the closest 10-fold cough result here is AUROC 0.819 $\pm$ 0.029. The comparison is useful because it separates two evidence targets. DNDT/DNDF achieves higher internal cough discrimination; this manuscript contributes a protocol-aware reliability layer that places internal, temporal, external, and metadata-control evidence in one decision-support framework.

The external and temporal rows ask different questions. DNDT/DNDF reports limited Coswara--COUGHVID cross-dataset AUCs near 0.53--0.57 \citep{islam2026dndf}, which is consistent with the COUGHVID transfer AUROC 0.543 reported here. Drift-adaptive work reports development-to-postdevelopment degradation in chronological cough experiments \citep{ganitidis2025drift}, which is directionally consistent with the temporal drop in this study. Symptom-assisted methods report high metrics when symptoms are intentionally included as predictive information \citep{lin2025sympcoughnet}; our symptoms-only AUROC 0.932 explains why those systems and audio-only systems should not be compared without stating the information available to the model.

\begin{table}[t]
\centering
\caption{Protocol-aware metric comparison against selected prior studies. Close comparisons are still not exact unless modalities, split units, and feature sets match; apples-to-oranges rows are included only to clarify why raw leaderboard comparisons are misleading.}
\label{tab:metric-comparison}
\scriptsize
\setlength{\tabcolsep}{2.5pt}
\begin{tabularx}{\linewidth}{p{0.13\linewidth} p{0.28\linewidth} p{0.27\linewidth} X}
\toprule
Mode & Prior evidence & Closest evidence here & Takeaway \\
\midrule
Internal cough & DNDF reports Coswara cough AUC 0.92 \citep{islam2026dndf}. & Cough 10-fold AUROC $0.819 \pm 0.029$; final source-domain fusion AUROC 0.897. & Internal cough accuracy and validation breadth are separate evidence targets. \\
Temporal & Drift-adaptive Coswara cough AUROC decreases from 0.668 to 0.597 \citep{ganitidis2025drift}. & Strict early-to-late Coswara AUROC 0.698 and AUPRC 0.896. & Both evidence sets support chronological testing. \\
External cough & DNDF reports Coswara-to-COUGHVID AUC 0.53 and reverse AUC 0.57 \citep{islam2026dndf}. & Model-matched cough-only AUROC changes from Coswara internal 0.849-0.868 to COUGHVID transfer 0.523-0.543. & Cross-dataset cough transfer is weak despite stronger internal cough rows. \\
Symptom input & SympCoughNet reports 89.30\% accuracy, 94.74\% AUROC, and 91.62\% PR with symptom-assisted input \citep{lin2025sympcoughnet}. & Symptoms-only AUROC 0.932; final audio/fusion internal AUROC 0.897. & Symptom-assisted prediction and symptom-confounding control answer different questions. \\
\bottomrule
\end{tabularx}
\end{table}

\section{Discussion}

The results support three expert-system lessons. First, a respiratory-audio model can look strong under a participant-disjoint source-domain protocol. The final Coswara fusion AUROC of 0.897 is not trivial and shows that the model bank can exploit reproducible signal inside the source dataset. Second, source-domain strength does not imply future-period or external-dataset reliability. The strict temporal AUROC of 0.698 and the model-matched cough-only internal-to-external drop from 0.849-0.868 to 0.523-0.543 show that the decision rule is sensitive to time and dataset identity. Third, metadata-only performance is strong enough to make causal audio claims unsafe without additional controls.

The comparison with DNDT/DNDF is especially important for ESWA readers because it clarifies the paper's role. DNDT/DNDF advances the internal cough-classification architecture. The present manuscript advances the validation layer around such systems: it asks whether an internally plausible expert system continues to look reliable when modality scope, chronological testing, external transfer, and metadata-only controls are made explicit. The DNDT/DNDF cross-dataset AUCs and our COUGHVID transfer result point in the same direction: COVID cough models can be substantially dataset-dependent.

The modality scope must remain precise. Coswara supports cough, breath, and speech fusion because those modalities are available in the source dataset. COUGHVID supports an independent cough-only test. A weak COUGHVID result cannot prove that multimodal fusion is useless; it proves that the selected Coswara cough decision rule does not transfer well to a different cough corpus. Conversely, a strong Coswara fusion result cannot prove external robustness without a multimodal external dataset.

For decision-support deployment, these findings argue for prospective validation, calibration monitoring, standardized recording protocols, external multimodal cohorts, and routine metadata audits. An expert system should report the validation tier under which its operating point was selected. It should also expose whether symptoms or protocol variables can predict labels independently of audio, because those variables may encode triage behavior, testing access, or collection campaigns rather than disease acoustics.

\subsection{Practical implications for expert-system design}

The framework changes how a respiratory-audio system should be evaluated before use. A model-development report should include a validation card with the dataset, modality set, participant-separation unit, calendar split, external dataset compatibility, prevalence, threshold-selection rule, and metadata-only baseline. A high AUROC under one row of that card is not sufficient for deployment-oriented language. This is particularly important for systems that mix clinical symptoms with audio: symptoms may improve prediction, but the resulting system is no longer evidence for an audio biomarker unless symptom-only and demographic/protocol-only controls are reported.

The same structure can be reused beyond COVID-19. Any respiratory-audio expert system for screening or triage should separate source-domain discrimination from external portability, identify whether the external dataset contains the same modalities, and test whether non-audio variables recover labels. The proposed framework is therefore a general validation design for respiratory-audio expert systems under distribution shift.

\section{Limitations and Future Validation}

The main threat to validity is not a single technical bug but the observational nature of the evidence. Coswara and COUGHVID differ in recruitment, prompts, labels, prevalence, recording environments, and available modalities. Participant-disjoint validation reduces identity leakage, but it does not remove site, device, geography, recruitment, symptom, or calendar effects. COUGHVID is cough-only, so it is a legitimate external cough-transfer endpoint but not an external validation of breath, speech, or multimodal fusion.

The primary validation ladder also lacks structural confidence intervals for every selected endpoint. The paper-comparable cross-validation rows report fold-level variability, but the final participant-split, temporal, and external-transfer rows are point estimates. Future versions should add DeLong, bootstrap, or paired-resampling confidence intervals so that temporal and external drops can be separated from small pairwise fluctuations. Until then, the manuscript uses the large direction of degradation as reliability evidence and avoids ranking close model variants by small metric differences.

The strict temporal split is intentionally severe, but it is still retrospective. It changes prevalence and sample composition, making AUPRC and thresholded metrics dependent on the test-period distribution. Metadata and inverse-propensity analyses show measured shortcut risk, but they cannot prove every causal pathway or remove unmeasured confounding. The cap-10 IPW row improves measured balance while reducing the effective sample size to 130.4, so the retained AUROC 0.780 should be read as qualified residual signal rather than a definitive causal audio effect. The model bank includes several practical expert-system components, but it is not an exhaustive search over all possible transformers, adaptation methods, event-segmentation methods, or calibration pipelines. Therefore, the claim is not that no stronger classifier exists. The claim is that an internally plausible respiratory-audio expert system must be evaluated under temporal, external, and metadata-control stress before it is treated as reliable decision support.

Some development-bank rows use different internal protocols and should not be treated as controlled ablations unless the table explicitly states a matched comparison. Thresholded metrics depend on validation threshold selection, and external calibration cannot be assumed under prevalence shift. The analysis does not establish clinical utility, prospective safety, or calibrated individual risk under deployment conditions.

These threats also define the strongest future work. The next validation step is not another internal architecture search alone; it is a prospective or well-curated external respiratory-audio cohort with matched cough, breath, and speech prompts, standardized acquisition metadata, prespecified thresholds, subgroup calibration, and a target-population screening endpoint. A stronger neural model would be scientifically valuable only if it improves performance at those stricter validation tiers.

\section{Conclusion}

This paper presents a protocol-aware expert-system framework for evaluating COVID-19 respiratory-audio screening models under temporal shift, external dataset shift, and metadata confounding. The final Coswara fusion model achieved strong participant-disjoint source-domain performance, but performance declined under stricter temporal validation and direct COUGHVID cough transfer. Metadata-only and symptoms-only controls were highly predictive, showing that observational respiratory-audio labels can be strongly entangled with non-audio variables. The scientific contribution is therefore a decision-support validation protocol and evidence package, not a deployment-ready diagnostic system or a raw SOTA claim.

\section*{Data and Code Availability}

Coswara and COUGHVID should be obtained from their original maintainers under the applicable dataset terms. The manuscript reports aggregate metrics, derived tables, and original study figures generated by the reproducible analysis workflow. Dataset redistribution constraints should be respected by releasing code, aggregate result tables, and figure-generation artifacts without redistributing restricted raw audio.

\bibliographystyle{elsarticle-num}
\bibliography{references}

\end{document}
